\documentclass[a4paper]{article}     
\usepackage{amsfonts}
\usepackage{amsmath}
\usepackage{amssymb}
\usepackage{graphicx}

\begin{document}

\noindent \large Auteur: Mohamed Shafik \\
\noindent \large Studierichting: Informatica\\
\noindent \large Oefeningenreeks: 6B nr 11\\

\medskip

\normalsize

Bepaal de eerste term en het verschil van een rekenkundige rij als men weet dat $s_n = n(n+2)$.\\

$x_1 = s_1 = 1 \cdot (1+2) = 3$\\

$x_2 = s_2 - s_1 = 2 \cdot 4 - 3 = 5$\\

$x_3 = s_3 - s_2 = 3 \cdot 5 - 2 \cdot 4 = 15 - 8 = 7$\\

$v = x_2 - x_1 = 5 - 3 = 2$\\

\begin{thebibliography}{999}
\end{thebibliography}
\end{document}
